\section{AMC 10 10B 2005}
        \begin{enumerate}
        	\item (\textit{AMC 10 10B 2005}) A scout troop buys $1000$ candy bars at a price of five for $$2$. They sell all the candy bars at a price of two for $$1$. What was the profit, in dollars?


 $\mathrm{(A)} 100 \qquad \mathrm{(B)} 200 \qquad \mathrm{(C)} 300 \qquad \mathrm{(D)} 400 \qquad \mathrm{(E)} 500$


 
 \textbf{Answer: }A



\textbf{Solution}


 \begin{align*} \mbox{Expenses} &= 1000 \cdot \frac25 = 400 \\ \mbox{Revenue}  &= 1000 \cdot \frac12 = 500 \\ \mbox{Profit}   &= \mbox{Revenue} - \mbox{Expenses} = 500-400 = \boxed{\textbf{(A) }100}. \end{align*}
Note: Revenue is a gain.

 



\textbf{Solution 2 (Faster)}

Note that the troop buys $10$ candy bars at a price of $4$ dollars and sells $10$ bars at a price of $5$ dollars. So the troop gains $1$ dollar for every $10$ bars. So therefore we divide $1000 \div 10 = 100$. So our answer is $\boxed{\textbf{(A) }100}$.~HyperVoid

 


	\item (\textit{AMC 10 10B 2005}) A positive number $x$ has the property that $x\%$ of $x$ is $4$. What is $x$?


 $\mathrm{(A)} 2 \qquad \mathrm{(B)} 4 \qquad \mathrm{(C)} 10 \qquad \mathrm{(D)} 20 \qquad \mathrm{(E)} 40$


 
 \textbf{Answer: }D



\textbf{Solution 1}

Since $x\%$ means $0.01x$, the statement "$x\% \text{ of } x \text{ is 4}$" can be rewritten as "$0.01x \cdot x = 4$":

 $0.01x \cdot x=4 \Rightarrow x^2 = 400 \Rightarrow x = \boxed{\textbf{(D) }20}.$

 



\textbf{Solution 2}

Try the answer choices one by one. Upon examination, it is quite obvious that the answer is $\boxed{\textbf{(D) }20}.$ Very fast.

 
Solution by franzliszt

 


	\item (\textit{AMC 10 10B 2005}) A gallon of paint is used to paint a room. One third of the paint is used on the first day. On the second day, one third of the remaining paint is used. What fraction of the original amount of paint is available to use on the third day?


 $\mathrm{(A)} \frac{1}{10} \qquad \mathrm{(B)} \frac{1}{9} \qquad \mathrm{(C)} \frac{1}{3} \qquad \mathrm{(D)} \frac{4}{9} \qquad \mathrm{(E)} \frac{5}{9}$


 
 \textbf{Answer: }D



\textbf{Solution}

After the first day, there is $1-\left(\dfrac{1}{3}\cdot1\right)=\frac{2}{3}$ gallons left. After the second day, there is a total of $\dfrac{2}{3}-\left(\dfrac{1}{3}\times\dfrac{2}{3}\right)=\dfrac{2}{3}-\dfrac{2}{9}=\dfrac{4}{9}$. Therefore, the fraction of the original amount of paint that is left is $\dfrac{\dfrac{4}{9}}{1}=\boxed{\textbf{(D) }\frac{4}{9}}$.

 Another way to do this is just to simply find the gain everyday and subtract from the remaining you had before.

 


	\item (\textit{AMC 10 10B 2005}) For real numbers $a$ and $b$, define $a \diamond b = \sqrt{a^2 + b^2}$. What is the value of


 $(5 \diamond 12) \diamond ((-12) \diamond (-5))$?


 $\mathrm{(A)} 0 \qquad \mathrm{(B)} \frac{17}{2} \qquad \mathrm{(C)} 13 \qquad \mathrm{(D)} 13\sqrt{2} \qquad \mathrm{(E)} 26$


 
 \textbf{Answer: }D



\textbf{Solution}


 \begin{align*} (5 \diamond 12) \diamond ((-12) \diamond (-5))&=(\sqrt{5^2+12^2}) \diamond (\sqrt{(-12)^2+(-5)^2})\\ &=(\sqrt{169})\diamond(\sqrt{169})\\ &=13\diamond13\\  &=\sqrt{13^2+13^2}\\  &=\sqrt{338}\\  &=\boxed{\mathrm{(D)\,13\sqrt{2}}}\\ \end{align*}
Note that the negative signs did not matter and any number squared times two is that number times the square root of 2.

 



\textbf{Solution 2 (Quick)}

Notice that the formula given is for the hypotenuse of any right triangle. As stated in the first solution, the negative signs did not matter. Therefore, we can think of this in terms of right triangles. Knowing the 5-12-13 Pythagorean triple, we find it to be simplified to
 \[(13 \diamond 13)\]Then, knowing the 45-45-90 degree triangle ratio, we find the answer to be 
 \[{(D)}\]~Thyauraisoverwhelming

 


	\item (\textit{AMC 10 10B 2005}) Brianna is using part of the money she earned on her weekend job to buy several equally-priced CDs. She used one fifth of her money to buy one third of the CDs. What fraction of her money will she have left after she buys all the CDs?


 $\mathrm{(A)} \frac{1}{5} \qquad \mathrm{(B)} \frac{1}{3} \qquad \mathrm{(C)} \frac{2}{5} \qquad \mathrm{(D)} \frac{2}{3} \qquad \mathrm{(E)} \frac{4}{5}$


 
 \textbf{Answer: }C



\textbf{Solution}

Let $m =$ Brianna's money.  We have $\frac15 m = \frac13 (\mbox{CDs}) \Rightarrow \frac35 m = (\mbox{CDs})$.  Thus, the money left over is $m-\frac35m = \frac25m$, so the answer is $\boxed{\textbf{(C) }\frac{2}{5}}$.

 This was just a simple manipulation of the equation. No solving was needed!

 


	\item (\textit{AMC 10 10B 2005}) At the beginning of the school year, Lisa's goal was to earn an A on at least $80\%$ of her $50$ quizzes for the year. She earned an A on $22$ of the first $30$ quizzes. If she is to achieve her goal, on at most how many of the remaining quizzes can she earn a grade lower than an A?


 $\mathrm{(A)} 1 \qquad \mathrm{(B)} 2 \qquad \mathrm{(C)} 3 \qquad \mathrm{(D)} 4 \qquad \mathrm{(E)} 5$


 
 \textbf{Answer: }B



\textbf{Solution}

Lisa's goal was to get an $A$ on $80\% \cdot 50 = 40$ quizzes.  She already has $A$'s on $22$ quizzes, so she needs to get $A$'s on $40-22=18$ more.  There are $50-30=20$ quizzes left, so she can afford to get less than an $A$ on $20-18=\boxed{\textbf{(B) }2}$ of them.

 Here, only the $A$'s matter... No complicated stuff!

 


	\item (\textit{AMC 10 10B 2005}) A circle is inscribed in a square, then a square is inscribed in this circle, and finally, a circle is inscribed in this square. What is the ratio of the area of the smaller circle to the area of the larger square?


 $\mathrm{(A)} \frac{\pi}{16} \qquad \mathrm{(B)} \frac{\pi}{8} \qquad \mathrm{(C)} \frac{3\pi}{16} \qquad \mathrm{(D)} \frac{\pi}{4} \qquad \mathrm{(E)} \frac{\pi}{2}$


 
 \textbf{Answer: }B



\textbf{Solution 1}

Let the side of the largest square be $x$. It follows that the diameter of the inscribed circle is also $x$. Therefore, the square's diagonal inscribed in the circle is $x$. The side length of the smaller square is $\dfrac{x}{\sqrt{2}}=\dfrac{x\sqrt{2}}{2}$. Similarly, the diameter of the smaller inscribed circle is $\dfrac{x\sqrt{2}}{2}$. Hence, its radius is $\dfrac{x\sqrt{2}}{4}$. The area of this circle is $\left(\dfrac{x\sqrt{2}}{4}\right)^2\pi=\dfrac{2\pi x^2}{16}=\dfrac{x^2\pi}{8}$, and the area of the largest square is $x^2$. The ratio of the areas is $\dfrac{\dfrac{x^2\pi}{8}}{x^2}=\frac{\cancel{x^2}\pi}{8}\cdot\frac{1}{\cancel{x^2}}=\boxed{\textbf{(B) }\frac{\pi}{8}}$.

 



\textbf{Solution 2}

Let the radius of the smallest circle be $r$. Then the side length of the smaller square is $2r$. The radius of the larger circle is half the length of the diagonal of the smaller square, so it is $\sqrt{2}r$. Hence the largest square has sides of length $2\sqrt{2}r$. The ratio of the area of the smallest circle to the area of the largest square is therefore $\frac{\pi r^2}{\left(2\sqrt{2}r\right)^2} =\boxed{\textbf{(B) }\frac{\pi}{8}}.$

 
                    \begin{figure}[H]
                        \centering
                        \includegraphics[width=0.5\textwidth]{Solution/AMC_10_10B_2005_7_Solution_Figure_1.png}
                    \end{figure}
                    Tip: When facing a geometry problem, it is very helpful to draw a diagram.

 


	\item (\textit{AMC 10 10B 2005}) An $8$-foot by $10$-foot floor is tiled with square tiles of size $1$ foot by $1$ foot. Each tile has a pattern consisting of four white quarter circles of radius $1/2$ foot centered at each corner of the tile. The remaining portion of the tile is shaded. How many square feet of the floor are shaded?

                    \begin{figure}[H]
                        \centering
                        \includegraphics[width=0.5\textwidth]{Problem/AMC_10_10B_2005_8_Problem_Figure_1.png}
                    \end{figure}
                    $\textrm{(A)}\ 80-20\pi \qquad \textrm{(B)}\ 60-10\pi \qquad \textrm{(C)}\ 80-10\pi \qquad \textrm{(D)}\ 60+10\pi \qquad \textrm{(E)}\ 80+10\pi$


 
 \textbf{Answer: }A



\textbf{Solution}

There are $80$ tiles.  Each tile has $[\mbox{square} - 4 \cdot (\mbox{quarter circle})]$ shaded.  Thus:

 
 \begin{align*} \mbox{shaded area} &= 80 \left( 1 - 4 \cdot \dfrac{1}{4} \cdot \pi \cdot \left(\dfrac{1}{2}\right)^2\right) \\ &= 80\left(1-\dfrac{1}{4}\pi\right) \\ &= \boxed{\textbf{(A) }80-20\pi}. \end{align*}
4 quarters of a circle is a circle so that may save you 0.5 seconds :)

 


	\item (\textit{AMC 10 10B 2005}) One fair die has faces $1$, $1$, $2$, $2$, $3$, $3$ and another has faces $4$, $4$, $5$, $5$, $6$, $6$. The dice are rolled and the numbers on the top faces are added. What is the probability that the sum will be odd?


 $\mathrm{(A)} \frac{1}{3} \qquad \mathrm{(B)} \frac{4}{9} \qquad \mathrm{(C)} \frac{1}{2} \qquad \mathrm{(D)} \frac{5}{9} \qquad \mathrm{(E)} \frac{2}{3}$


 
 \textbf{Answer: }D



\textbf{Solution}

In order to obtain an odd sum, exactly one out of the two dice must have an odd number. We can easily find the total probability using casework.

 \textbf{Case 1}: The first die is odd and the second die is even.

 The probability of this happening is $\dfrac{4}{6}\times\dfrac{4}{6}=\dfrac{16}{36}=\dfrac{4}{9}$

 \textbf{Case 2}: The first die is even and the second die is odd.

 The probability of this happening is $\dfrac{2}{6}\times\dfrac{2}{6}=\dfrac{4}{36}=\dfrac{1}{9}$

 Adding these two probabilities will give us our final answer. $\dfrac{4}{9}+\dfrac{1}{9}=\boxed{\textbf{(D) }\dfrac{5}{9}}$

 If you run out of time, you can see it is obviously greater than $\frac{1}{2}$ so it narrows guesses.

 



\textbf{Solution 2 (Complementary Counting)}

We see can see that there are $2$ ways to get an even number, and other ways get us an odd. Therefore, if we subtract the $P(\text{even})$, we will get $P(\text{odd})$. We can get an even either by getting $2$ evens or $2$ odds. Both cases give $\frac{2}{9}$, so $P(\text{odd})$ is $1-2 \cdot \frac{2}{9}=\boxed{\textbf{(D) }\frac{5}{9}}$

 


	\item (\textit{AMC 10 10B 2005}) In $\triangle ABC$, we have $AC = BC = 7$ and $AB = 2$. Suppose that $D$ is a point on line $AB$ such that $B$ lies between $A$ and $D$ and $CD = 8$. What is $BD$?


 $\mathrm{(A)} 3 \qquad \mathrm{(B)} 2\sqrt{3} \qquad \mathrm{(C)} 4 \qquad \mathrm{(D)} 5 \qquad \mathrm{(E)} 4\sqrt{2}$


 
 \textbf{Answer: }A



\textbf{Solution 1}

Draw height $CH$ (Perpendicular line from point C to line AD). We have that $BH=1$. By the \href{https://artofproblemsolving.com/wiki/index.php?title=Pythagorean\_Theorem}{Pythagorean Theorem}, $CH=\sqrt{48}$. Since $CD=8$, $HD=\sqrt{8^2-48}=\sqrt{16}=4$, and $BD=HD-1$, so $BD=\boxed{\textbf{(A) }3}$.

 



\textbf{Solution 2 (Trig)}

After drawing out a diagram, let $\angle{ABC}=\theta$. By the Law of Cosines, $7^2=2^2+7^2-2(7)(2)\cos{\theta} \rightarrow 0=4-28\cos{\theta} \rightarrow \cos{\theta}=\frac{1}{7}$. In $\triangle CBD$, we have $\angle{CBD}=(180-\theta)$, and using the identity $\cos(180-\theta)=-\cos{\theta}$ and Law of Cosines one more time: $8^2=7^2+x^2-2(7)(x)\left( \frac{-1}{7} \right) \rightarrow 64=49+x^2+2x \rightarrow x^2+2x-15=0$. The only positive value for $x$ is $3$, which gives the length of $\overline{BD}$. Thus the answer is $\boxed{\textbf{(A) }3}$.

 



\textbf{Solution 3 (Stewart's Theorem)}

Let $BD=k$. Then, by \href{https://artofproblemsolving.com/wiki/index.php?title=Stewart\%27s\_Theorem}{Stewart's Theorem}, 

 $2k(2+k)+7^2(2+k)=7^2k+8^2\cdot 2 \implies k^2+2k-15=0 \implies k=\boxed{\textbf{(A) }3}$

 


	\item (\textit{AMC 10 10B 2005}) The first term of a sequence is $2005$. Each succeeding term is the sum of the cubes of the digits of the previous term. What is the ${2005}^{\text{th}}$ term of the sequence?


 $\mathrm{(A)} 29 \qquad \mathrm{(B)} 55 \qquad \mathrm{(C)} 85 \qquad \mathrm{(D)} 133 \qquad \mathrm{(E)} 250$


 
 \textbf{Answer: }E



\textbf{Solution}

Performing this operation several times yields the results of $133$ for the second term, $55$ for the third term, and $250$ for the fourth term. The sum of the cubes of the digits of $250$ equal $133$, a complete cycle. The cycle is, excluding the first term, the $2^{\text{nd}}$, $3^{\text{rd}}$, and $4^{\text{th}}$ terms will equal $133$, $55$, and $250$, following the fourth term. Any term number that is equivalent to $1\ (\text{mod}\ 3)$ will produce a result of $250$. It just so happens that $2005\equiv 1\ (\text{mod}\ 3)$, which leads us to the answer of $\boxed{\textbf{(E) } 250}$.

 


	\item (\textit{AMC 10 10B 2005}) Twelve fair dice are rolled. What is the probability that the product of the numbers on the top faces is prime?


 $\mathrm{(A)} \left(\frac{1}{12}\right)^{12} \qquad \mathrm{(B)} \left(\frac{1}{6}\right)^{12} \qquad \mathrm{(C)} 2\left(\frac{1}{6}\right)^{11} \qquad \mathrm{(D)} \frac{5}{2}\left(\frac{1}{6}\right)^{11} \qquad \mathrm{(E)} \left(\frac{1}{6}\right)^{10}$


 
 \textbf{Answer: }E



\textbf{Solution 1}

In order for the product of the numbers to be prime, $11$ of the dice have to be a $1$, and the other die has to be a prime number. There are $3$ prime numbers ($2$, $3$, and $5$), and there is only one $1$, and there are $\dbinom{12}{1}$ ways to choose which die will have the prime number, so the probability is $\dfrac{3}{6}\times\left(\dfrac{1}{6}\right)^{11}\times\dbinom{12}{1} = \dfrac{1}{2}\times\left(\dfrac{1}{6}\right)^{11}\times12=\left(\dfrac{1}{6}\right)^{11}\times6=\boxed{\textbf{(E) }\left(\dfrac{1}{6}\right)^{10}}$.

 



\textbf{Solution 2}

There are three cases where the product of the numbers is prime. One die will show $2$, $3$, or $5$ and each of the other $11$ dice will show a $1$. For each of these three cases, the number of ways to order the numbers is $\dbinom{12}{1}$ = $12$ . There are $6$ possible numbers for each of the $12$ dice, so the total number of permutations is $6^{12}$. The probability the product is prime is therefore $\frac{3\cdot 12}{6^{12}} = \frac{1}{6^{10}} =\boxed{\textbf{(E) }\left(\dfrac{1}{6}\right)^{10}}$.

 



\textbf{Solution 3}

The only way to get a product of that is a prime number is to roll all ones except for such prime, e.g: $11$ ones and $1$ two. So we seek the probability of rolling $11$ ones and $1$ prime number. The probability of rolling $11$ ones is $\frac{1}{6^{11}}$ and the probability of rolling a prime is $\frac{1}{2}$, giving us a probability of $\frac{1}{6^{11}}\cdot\frac{1}{2}$ of this outcome occuring. However, there are $\frac{12!}{11!\cdot{1!}}=12$ ways to arrange the ones and the prime. Multiplying the previous probability by $12$ gives us $\frac{1}{6^{11}}\cdot\frac{1}{2}\cdot{6}=\frac{1}{6^{10}}=\boxed{\textbf{(E) }\left(\dfrac{1}{6}\right)^{10}}.$

 -Benedict T (countmath1)

 



\textbf{Solution 4}

The only way to get a prime number product is to roll \(1\) on all \(11\) dices and roll a prime number on the remaining one. 

 There are 3 prime numbers less than or equal to 6: \(2, 3, 5\). For each of the 12 dice, we can choose one to roll the prime number on, while the other 11 show 1. So, there are \(12 \times 3 = 36\) favorable outcomes.

 Since each die has 6 faces and there are 12 dice, the total number of possible outcomes is \(6^{12}\).

 Therefore, the probability is 
 \[\frac{36}{6^{12}} = \left(\frac{1}{6}\right)^{10}.\]

 
 \[\boxed{\textbf{(E) } \left(\dfrac{1}{6}\right)^{10}}\]-LittleWavelet

 


	\item (\textit{AMC 10 10B 2005}) How many numbers between $1$ and $2005$ are integer multiples of $3$ or $4$ but not $12$?


 $\mathrm{(A)} 501 \qquad \mathrm{(B)} 668 \qquad \mathrm{(C)} 835 \qquad \mathrm{(D)} 1002 \qquad \mathrm{(E)} 1169$


 
 \textbf{Answer: }C



\textbf{Solution 1}

To find the multiples of $3$ or $4$ but not $12$, you need to find the number of multiples of $3$ and $4$, and then subtract twice the number of multiples of $12$, because you overcount and do not want to include them. The multiples of $3$ are $\frac{2005}{3} = 668\text{ }R1.$ The multiples of $4$ are $\frac{2005}{4} = 501 \text{ }R1$. The multiples of $12$ are $\frac{2005}{12} = 167\text{ }R1.$ So, the answer is $668+501-167-167 = \boxed{\textbf{(C) } 835}$

 



\textbf{Solution 2}

From $1$-$12$, the multiples of $3$ or $4$ but not $12$ are $3, 4, 6, 8,$and $9$, a total of five numbers. Since $\frac{5}{12}$ of positive integers are multiples of $3$ or $4$ but not $12$ from $1$-$12$, the answer is approximately $\frac{5}{12} \cdot 2005$ = $\boxed{\textbf{(C) }835}$

 



\textbf{Solution 3}

\begin{verbatim}
(Simpler version of solution 2)
\end{verbatim}

For every multiple of 12, there are a total of 7 multiples of 3 and 4 combined. However, we do not want to include the 2 multiples of 12, so we subtract 2 from 7, leaving us with 5.

 We divide 2005 by 12 to determine how many such multiples exist:
 \[\left\lfloor \frac{2005}{12} \right\rfloor = 167\]

 Multiplying that by 5, we get:
 \[167 \times 5 = 835\]

 So our answer is:
 \[\boxed{\textbf{(C) } 835}\]

 -LittleWavelet

 


	\item (\textit{AMC 10 10B 2005}) Equilateral $\triangle ABC$ has side length $2$, $M$ is the midpoint of $\overline{AC}$, and $C$ is the midpoint of $\overline{BD}$. What is the area of $\triangle CDM$?

                    \begin{figure}[H]
                        \centering
                        \includegraphics[width=0.5\textwidth]{Problem/AMC_10_10B_2005_14_Problem_Figure_1.png}
                    \end{figure}
                    $\textrm{(A)}\ \frac {\sqrt {2}}{2}\qquad \textrm{(B)}\ \frac {3}{4}\qquad \textrm{(C)}\ \frac {\sqrt {3}}{2}\qquad \textrm{(D)}\ 1\qquad \textrm{(E)}\ \sqrt {2}$


 
 \textbf{Answer: }C



\textbf{Solution 1 (trig)}

The area of a triangle can be given by $\frac12 ab \sin C$. $MC=1$ because it is the midpoint of a side, and $CD=2$ because it is the same length as $BC$. Each angle of an equilateral triangle is $60^\circ$ so $\angle MCD = 120^\circ$. The area is $\frac12 (1)(2) \sin  120^\circ = \boxed{\textbf{(C)}\ \frac{\sqrt{3}}{2}}$.Note: Even if you don't know the value of $\sin 120^\circ$, you can use the fact that  $\sin x = \sin (180^\circ - x)$, so $\sin 120^\circ = \sin 60^\circ$.You can easily calculate $\sin 60^\circ$ to be $\frac{\sqrt3}{2}$ using equilateral triangles.

 



\textbf{Solution 2}

In order to calculate the area of $\triangle CDM$, we can use the formula $A=\dfrac{1}{2}bh$, where $\overline{CD}$ is the base. We already know that $\overline{CD}=2$, so the formula becomes $A=h$. We can drop verticals down from $A$ and $M$ to points $E$ and $F$, respectively. We can see that $\triangle AEC \sim \triangle MFC$. Now, we establish the relationship that $\dfrac{AE}{MF}=\dfrac{AC}{MC}$. We are given that $\overline{AC}=2$, and $M$ is the midpoint of $\overline{AC}$, so $\overline{MC}=1$. Because $\triangle AEB$ is a $30-60-90$ triangle and the ratio of the sides opposite the angles are $1-\sqrt{3}-2$ $\overline{AE}$ is $\sqrt{3}$. Plugging those numbers in, we have $\dfrac{\sqrt{3}}{MF}=\dfrac{2}{1}$. Cross-multiplying, we see that $2\times\overline{MF}=\sqrt{3}\times1\implies \overline{MF}=\dfrac{\sqrt{3}}{2}$ Since $\overline{MF}$ is the height $\triangle CDM$, the area is $\boxed{\textbf{(C) }\frac{\sqrt{3}}{2}}$.

 



\textbf{Solution 3}

Draw a line from $M$ to the midpoint of $\overline{BC}$. Call the midpoint of $\overline{BC}$ $P$. This is an equilateral triangle, since the two segments $\overline{PC}$ and $\overline{MC}$ are identical, and $\angle C$ is $60^{\circ}$. Using the \href{https://artofproblemsolving.com/wiki/index.php?title=Pythagorean\_Theorem}{Pythagorean Theorem}, point $M$ to $\overline{BC}$ is $\dfrac{\sqrt{3}}{2}$. Also, the length of $\overline{CD}$ is 2, since $C$ is the midpoint of $\overline{BD}$. So, our final equation is $\frac{\sqrt{3}}{2}\times2\over2$, which just leaves us with $\boxed{\textbf{(C) }\dfrac{\sqrt{3}}{2}}$.

 



\textbf{Solution 4}

Drop a vertical down from $M$ to $BC$. Let us call the point of intersection $X$ and the midpoint of $BC$, $Y$. We can observe that $\triangle AYC$ and $\bigtriangleup MXC$ are similar. By the \href{https://artofproblemsolving.com/wiki/index.php?title=Pythagorean\_theorem}{Pythagorean theorem}, $AY$ is $\sqrt3$.

 Since $AC:MC=2:1,$ we find $MX=\frac{\sqrt3}{2}.$ Because $C$ is the midpoint of $BD,$ and $BC=2,$ $CD=2.$ Using the area formula, $\frac{CD*MX}{2}=\boxed{\textbf{(C) }\dfrac{\sqrt{3}}{2}}.$

 



\textbf{Solution 5}

To find area we can just solve for ratios between the triangles, and then plug in the area of $\triangle ABC$. The problem statement gives us $C$ as the midpoint of $\overline{BD}$, so $BC=CD$. For the height, we can see that $M$ is the midpoint of $\overline{AC}$ which means that $AM=CM$, and in turn means that the height of $\triangle MCD$ is half of that of $\triangle ABC$. 

 As the bases are the same, the ratio of the areas is the same of the heights, so $[MCD] = \frac{1}{2} [ABC]$. The area of an equilateral triangle $\triangle ABC$ is 
 \[\frac{s^2 \sqrt{3} }{4} = \frac{2^2 \sqrt{3} }{4} = \sqrt{3}\] so as $[MCD]$ is half of that, our final answer is 
 \[\boxed{\textbf{(C) }\dfrac{\sqrt{3}}{2}}.\]

 


	\item (\textit{AMC 10 10B 2005}) An envelope contains eight bills: $2$ ones, $2$ fives, $2$ tens, and $2$ twenties. Two bills are drawn at random without replacement. What is the probability that their sum is $$20$ or more?


 $\mathrm{(A)} \frac{1}{4} \qquad \mathrm{(B)} \frac{2}{5} \qquad \mathrm{(C)} \frac{3}{7} \qquad \mathrm{(D)} \frac{1}{2} \qquad \mathrm{(E)} \frac{2}{3}$


 
 \textbf{Answer: }D



\textbf{Solution 1}

The only way to get a total of \$$20$ or more is if you pick a twenty and another bill, or if you pick both tens. There are a total of $\dbinom{8}{2}=\dfrac{8\times7}{2\times1}=28$ ways to choose $2$ bills out of $8$. There are $12$ ways to choose a twenty and some other non-twenty bill. There is $1$ way to choose both twenties, and also $1$ way to choose both tens. Adding these up, we find that there are a total of $14$ ways to attain a sum of $20$ or greater, so there is a total probability of $\dfrac{14}{28}=\boxed{\textbf{(D) }\frac{1}{2}}$.

 



\textbf{Solution 2}

Another way to do this problem is to use \href{https://artofproblemsolving.com/wiki/index.php?title=Complementary\_counting}{complementary counting}, i.e. how many ways that the sum is less than $20$. Now, you do not have to consider the $2$ twenties, so you have $6$ bills left. $\dbinom{6}{2} = \dfrac{6\times5}{2\times1} = 15$ ways. However, you counted the case when you have $2$ tens, so you need to subtract 1, and you get $14$. Finding the ways to get $20$ or higher, you subtract $14$ from $28$ and get $14$. So the answer is $\dfrac{14}{28} = \boxed{\textbf{(D) }\dfrac{1}{2}}$

 



\textbf{Solution 3}

There are two cases that work, namely getting at least $1$ twenty, or getting $2$ tens.

 Case $1$: $P(\text{Get at least one twenty}) = 1-P(\text{Do not get a single twenty})=1- \frac{\binom{6}{2}}{\binom{8}{2}}=\frac{28-15}{28}=\frac{13}{28}$

 Case $2$ : $P(\text{Get two tens}) = \frac{1}{\binom{8}{2}} = \frac{1}{28}$

 Summing up our cases, we have $\frac{13}{28}+\frac{1}{28}=\frac{14}{28}=\boxed{\textbf{(D) } \dfrac{1}{2}}$

 



\textbf{Solution 4}

Note that if a twenty is drawn, anything else that is drawn will create a total greater than $20$; The probability of a twenty being drawn first is $\frac{1}{4}.$ The same could be said for drawing anything, and then drawing a twenty. However, we can only draw something that isn't a twenty first (since we've already accounted for the probability of drawing two twenties). 

 The probability of drawing a non-twenty first, then a twenty second is $\frac{3}{4}\cdot\frac{2}{7}=\frac{3}{14}.$ Finally, we can draw two tens. The probability of this occuring is $\frac{1}{4}\cdot\frac{1}{7}=\frac{1}{28}.$ 

 Adding these three probabilities gives us $\frac{1}{4}+\frac{3}{14}+\frac{1}{28}=\boxed{\textbf{(D) } \dfrac{1}{2}}$

 -Benedict T (countmath1) (edited by AMC\_8)

 



\textbf{Solution 5 (Quick if no time)}

We see that there are $2$ of each bill. We can simplify the question to only drawing $1$ bill out of $4$, and trying to draw $\$ 10$. (Originally you need to draw $\$ 20$, but the bill is halved, so you need to draw only $\$ 10$.) We see that we have $1$ of each bill $\$1$, $\$5$, $\$10$, and $\$20$.Thus, we can easily see the solution is $\boxed{(D) \frac{1}{2}}$.

 


	\item (\textit{AMC 10 10B 2005}) The quadratic equation $x^2 + mx + n = 0$ has roots that are twice those of $x^2 + px + m = 0$, and none of $m$, $n$, and $p$ is zero. What is the value of $n/p$?


 $\mathrm{(A)} 1 \qquad \mathrm{(B)} 2 \qquad \mathrm{(C)} 4 \qquad \mathrm{(D)} 8 \qquad \mathrm{(E)} 16$


 
 \textbf{Answer: }D



\textbf{Solution 1}

Let $x^2 + px + m = 0$ have roots $a$ and $b$. Then 

 
 \[x^2 + px + m = (x-a)(x-b) = x^2 - (a+b)x + ab,\]

 so $p = -(a+b)$ and $m = ab$. Also, $x^2 + mx + n = 0$ has roots $2a$ and $2b$, so

 
 \[x^2 + mx + n = (x-2a)(x-2b) = x^2 - 2(a+b)x + 4ab,\]

 and $m = -2(a+b)$ and $n = 4ab$. Thus $\frac{n}{p} = \frac{4ab}{-(a+b)} = \frac{4m}{\frac{m}{2}} = \boxed{\textbf{(D) }8}$. 

 To test that this actually works, consider the quadratics $x^2 + 8x + 16 = 0,\ x^2 + 16x + 64 = 0$.

 



\textbf{Solution 2}

If the roots of $x^2 + mx + n = 0$ are $2a$ and $2b$ and the roots of $x^2 + px + m = 0$ are $a$ and $b$, then using Vieta's formulas,
 \[2a + 2b = -m\]
 \[a + b = -p\]
 \[2a(2b) = n\]
 \[a(b) = m\]Therefore, substituting the second equation into the first equation gives 
 \[m = 2(p)\]and substituting the fourth equation into the third equation gives
 \[n = 4(m)\]Therefore, $n = 8p$, so $\frac{n}{p}= \boxed{\textbf{(D) }8}$

 


	\item (\textit{AMC 10 10B 2005}) Suppose that $4^a = 5$, $5^b = 6$, $6^c = 7$, and $7^d = 8$. What is $a \cdot b \cdot c \cdot d$?


 $\mathrm{(A)} 1 \qquad \mathrm{(B)} \frac{3}{2} \qquad \mathrm{(C)} 2 \qquad \mathrm{(D)} \frac{5}{2} \qquad \mathrm{(E)} 3$


 
 \textbf{Answer: }B



\textbf{Solution}


 \begin{align*} 8&=7^d \\ 8&=\left(6^c\right)^d\\ 8&=\left(\left(5^b\right)^c\right)^d\\ 8&=\left(\left(\left(4^a\right)^b\right)^c\right)^d\\ 8&=4^{a\cdot b\cdot c\cdot d}\\ 2^3&=2^{2\cdot a\cdot b\cdot c\cdot d}\\ 3&=2\cdot a\cdot b\cdot c\cdot d\\ a\cdot b\cdot c\cdot d&=\boxed{\textbf{(B) }\dfrac{3}{2}}\\ \end{align*}


 



\textbf{Solution 2 (logarithms)}

We can write $a$ as $\log_4 5$, $b$ as $\log_5 6$, $c$ as $\log_6 7$, and $d$ as $\log_7 8$. 

 We know that $\log_b a$ can be rewritten as $\frac{\log a}{\log b}$, so we have:
 \begin{align*} a\cdot b \cdot c \cdot d &= \frac{\cancel{\log5}}{\log4}\cdot\frac{\cancel{\log6}}{\cancel{\log5}}\cdot\frac{\cancel{\log7}}{\cancel{\log6}}\cdot\frac{\log8}{\cancel{\log7}} \\ a\cdot b \cdot c \cdot d &= \frac{\log8}{\log4} \\ a\cdot b \cdot c \cdot d &= \frac{3\cancel{\log2}}{2\cancel{\log2}} \\ a\cdot b \cdot c \cdot d &= \boxed{\textbf{(B) }\frac{3}{2}} \\ \end{align*}


 



\textbf{Solution 3 (logarithm chain rule)}

As in Solution 2, we can write $a$ as $\log_4 5$, $b$ as $\log_56$, $c$ as $\log_67$, and $d$ as $\log_78$. $a\cdot b\cdot c\cdot d$ is equivalent to $(\log_4 5)\cdot (\log_5 6)\cdot (\log_6 7)\cdot (\log_7 8)$. By the logarithm chain rule, this is equivalent to $\log_4 8$, which evaluates to $\boxed{\textbf{(B) }\frac{3}{2}}$.

 


	\item (\textit{AMC 10 10B 2005}) All of David's telephone numbers have the form $555-abc-defg$, where $a$, $b$, $c$, $d$, $e$, $f$, and $g$ are distinct digits and in increasing order, and none is either $0$ or $1$. How many different telephone numbers can David have?


 $\mathrm{(A)} 1 \qquad \mathrm{(B)} 2 \qquad \mathrm{(C)} 7 \qquad \mathrm{(D)} 8 \qquad \mathrm{(E)} 9$


 
 \textbf{Answer: }D



\textbf{Solution}

The only digits available to use in the phone number are $2$, $3$, $4$, $5$, $6$, $7$, $8$, and $9$. There are only $7$ spots left among the $8$ numbers, so we need to find the number of ways to choose $7$ numbers from $8$. The answer is then $\dbinom{8}{7}=\dfrac{8!}{7!\,(8-7)!}=\boxed{\textbf{(D) } 8}$

 Alternatively, we could just choose $1$ out of the $8$ numbers not to be used. There are obviously $\boxed{8}$ ways to do so.

 Note: for each combination of 7 numbers, exactly 1 is in increasing order -Williamgolly

 


	\item (\textit{AMC 10 10B 2005}) On a certain math exam, $10\%$ of the students got $70$ points, $25\%$ got $80$ points, $20\%$ got $85$ points, $15\%$ got $90$ points, and the rest got $95$ points. What is the difference between the mean and the median score on this exam?


 $\mathrm{(A)}\ 0 \qquad \mathrm{(B)}\ 1 \qquad \mathrm{(C)}\ 2 \qquad \mathrm{(D)}\ 4 \qquad \mathrm{(E)}\ 5$


 
 \textbf{Answer: }B



\textbf{Solution}

To begin, we see that the remaining $30\%$ of the students got $95$ points. Assume that there are $20$ students; we see that $2$ students got $70$ points, $5$ students got $80$ points, $4$ students got $85$ points, $3$ students got $90$ points, and $6$ students got $95$ points. The median is $85$, since the $10^{\text{th}}$ and $11^{\text{th}}$ terms are both $85$. The mean is $\dfrac{70\,(2)+80\,(5)+85\,(4)+90\,(3)+95\,(6)}{20}=\dfrac{1720}{20}=86$. The difference between the mean and median, therefore, is $\boxed{\textbf{(B)}\ 1}$.

 



\textbf{Solution 2}

The remaining $30\%$ of the students got $95$ points.The mean is equal to $10\%\cdot70 + 25\%\cdot80 + 20\%\cdot85 + 15\%\cdot90 + 30\%\cdot95 = 86$.The score greater than or equal to $50\%$ of other scores is the median. Since $35\%$ scored $80$ or lower and the next $20\%$ scored $85$, the median is $85$. The difference between the mean and the median is $\boxed{\textbf{(B)}\ 1}$.

 


	\item (\textit{AMC 10 10B 2005}) What is the average (mean) of all $5$-digit numbers that can be formed by using each of the digits $1$, $3$, $5$, $7$, and $8$ exactly once?


 $\mathrm{(A)} 48000 \qquad \mathrm{(B)} 49999.5 \qquad \mathrm{(C)} 53332.8 \qquad \mathrm{(D)} 55555 \qquad \mathrm{(E)} 56432.8$


 
 \textbf{Answer: }C



\textbf{Solution 1 (Expectation)}

The average of all valid numbers is simply the expected value of a randomly chosen valid number. In other words, it is $E[\text{ten thousands digit}]\cdot 10^4+E[\text{thousands digit}]\cdot 10^3+\cdots+E[\text{ones digit}]\cdot 1$. However, every digit is equally likely to be chosen for every place, so we can simplify this to $\frac{1+3+5+7+8}{5}\cdot(10^4+10^3+10^2+10^1+1)$ $=4.8\cdot11111 = \boxed{(C) 53332.8}$.

 



\textbf{Solution 2}

We first look at how many times each number will appear in each slot. If we fix a number in a slot, then there are $4! = 24$ ways to arrange the other numbers, so each number appears in each spot $24$ times. Therefore, the sum of all such numbers is $24 \times (1+3+5+7+8) \times (11111) = 24 \times 24 \times 11111 = 6399936.$ Since there are $5! = 120$ such numbers, we divide $6399936 \div 120$ to get $\boxed{\textbf{(C) }53332.8}$

 



\textbf{Solution 3}

We can first solve for the mean for the digits $1, 3, 5, 7,$and $9$ since each is $2$ away from each other. The mean of the numbers than can be solved using these digits is $55555$. The total amount of numbers that can be formed using these digits is $5! =120$. The sum of these numbers is $55555(120) = 6666600$. Now we can find out the total value that was gained by replacing the $8$ with a $9$. We can start how be calculating the gain when the $8$ was in the ones digit. Since there are $4! = 24$ numbers with the $8$ in the ones digit and $1$ was gain from each of them, $24$ is the number gained. Then, we repeat this with the tens, hundreds, thousands, and ten thousands place, leading to a total of $24+240+2400+24000+240000=266664$ as the total amount that was gained. Subtract this amount from the sum of the digits using the $9$ instead of the $8$ to get $6666600-266664=6399936$. Finally, we divide this by $120$ to get the average. $\frac{6399936}{120}= \boxed{\textbf{(C) }53332.8}$

 


	\item (\textit{AMC 10 10B 2005}) Forty slips are placed into a hat, each bearing a number $1$, $2$, $3$, $4$, $5$, $6$, $7$, $8$, $9$, or $10$, with each number entered on four slips. Four slips are drawn from the hat at random and without replacement. Let $p$ be the probability that all four slips bear the same number. Let $q$ be the probability that two of the slips bear a number $a$ and the other two bear a number $b \neq a$. What is the value of $\frac{q}{p}$?


 $\mathrm{(A)} 162 \qquad \mathrm{(B)} 180 \qquad \mathrm{(C)} 324 \qquad \mathrm{(D)} 360 \qquad \mathrm{(E)} 720$


 
 \textbf{Answer: }A



\textbf{Solution 1 (where the order of drawing slips matters)}

There are $10$ ways to determine which number to pick. There are $4!$ ways to then draw those four slips with that number, and $40 \cdot 39 \cdot 38 \cdot 37$ total ways to draw four slips. Thus $p = \frac{10\cdot 4!}{40 \cdot 39 \cdot 38 \cdot 37}$. 

 There are ${10 \choose 2} = 45$ ways to determine which two numbers to pick for the second probability. There are ${4 \choose 2} = 6$ ways to arrange the order which we draw the non-equal slips, and in each order there are $4 \times 3 \times 4 \times 3$ ways to pick the slips, so $q = \frac{45 \cdot 6 \cdot 4^2 \cdot 3^2}{40 \cdot 39 \cdot 38 \cdot 37}$. 

 Hence, the answer is $\frac{q}{p} = \frac{45 \cdot 6 \cdot 4^2 \cdot 3^2}{10\cdot 4!} = \boxed{\textbf{(A) }162}$.

 



\textbf{Solution 2 (where the order does not matter)}

For probability $p$, there are $\binom{10}{1}=10$ ways to choose the number you want to show up $4$ times.

 Hence, the probability is $\frac{10}{\binom{40}{4}}$.

 For probability $q$, there are $\binom{10}{2}=45$ ways to choose the $2$ numbers you want to show up twice. There are $\binom{4}{2}\cdot\binom{4}{2}$ ways to pick which slips you want out of the $4$ of each.

 Hence, the probability is $\frac{45\cdot6\cdot6}{\binom{40}{4}}$

 Hence, $\frac{q}{p}=\frac{45\cdot6\cdot6}{10}=9\cdot6\cdot3=\boxed{\textbf{(A) }162}$.

 


	\item (\textit{AMC 10 10B 2005}) For how many positive integers $n$ less than or equal to $24$ is $n!$ evenly divisible by $1 + 2 + \ldots + n$?


 $\mathrm{(A)} 8 \qquad \mathrm{(B)} 12 \qquad \mathrm{(C)} 16 \qquad \mathrm{(D)} 17 \qquad \mathrm{(E)} 21$


 
 \textbf{Answer: }C



\textbf{Solution}

Since $1 + 2 + \cdots + n = \frac{n(n+1)}{2}$, the condition is equivalent to having an integer value for $\frac{n!} {\frac{n(n+1)}{2}}$. This reduces, when $n\ge 1$, to having an integer value for $\frac{2(n-1)!}{n+1}$. This fraction is an integer unless $n+1$ is an odd prime. There are $8$ odd primes less than or equal to $24$, so there 

 are $24 - 8 =  \boxed{\textbf{(C) }16}$ numbers less than or equal to $24$ that satisfy the condition.

 


	\item (\textit{AMC 10 10B 2005}) In trapezoid $ABCD$ we have $\overline{AB}$ parallel to $\overline{DC}$, $E$ as the midpoint of $\overline{BC}$, and $F$ as the midpoint of $\overline{DA}$. The area of $ABEF$ is twice the area of $FECD$. What is $AB/DC$?


 $\mathrm{(A)} 2 \qquad \mathrm{(B)} 3 \qquad \mathrm{(C)} 5 \qquad \mathrm{(D)} 6 \qquad \mathrm{(E)} 8$


 
 \textbf{Answer: }C



\textbf{Solution 1}

Since the heights of both trapezoids are equal, and the area of $ABEF$ is twice the area of $FECD$, 

 $\frac{AB+EF}{2}=2\left(\frac{DC+EF}{2}\right)$.

 $\frac{AB+EF}{2}=DC+EF$, so

 $AB+EF=2DC+2EF$.

 $EF$ is exactly halfway between $AB$ and $DC$, so $EF=\frac{AB+DC}{2}$. 

 $AB+\frac{AB+DC}{2}=2DC+AB+DC$, so

 $\frac{3}{2}AB+\frac{1}{2}DC=3DC+AB$, and 

 $\frac{1}{2}AB=\frac{5}{2}DC$.

 $\frac{AB}{DC} = \boxed{\textbf{(C) }5}$.

 



\textbf{Solution 2}

Mark $DC=z$, $AB=x$, and $FE=y.$Note that the heights of trapezoids $ABEF$ & $FECD$ are the same. Mark the height to be $h$.

 Then, we have that $\frac{x+y}{2}\cdot h=2(\frac{y+z}{2} \cdot h)$.

 From this, we get that $x=2z+y$.

 We also get that $\frac{x+z}{2} \cdot 2h= 3(\frac{y+z}{2} \cdot h)$.

 Simplifying, we get that $2x=z+3y$

 Notice that we want $\frac{AB}{DC}=\frac{x}{z}$.

 Dividing the first equation by $z$, we get that $\frac{x}{z}=2+\frac{y}{z}\implies 3(\frac{x}{z})=6+3(\frac{y}{z})$.

 Dividing the second equation by $z$, we get that $2(\frac{x}{z})=1+3(\frac{y}{z})$.

 Now, when we subtract the top equation from the bottom, we get that $\frac{x}{z}=\boxed{\textbf{(C) }5}$

 



\textbf{Solution 3}

Since the bases of the trapezoids along with the height are the same, the only thing that matters is the second base. Denote the length of the bigger trapezoid $x$. The area of the smaller trapezoid is $A$ = $h \cdot \frac {b_1 + b_2}{2}$. The area of the larger trapezoid is $A$ = $h \cdot \frac {b_1 + x}{2}$. Since this problem asks for proportions, assume that $b_1 = 1$ and $b_2 = 2$. 

 The smaller trapezoid has area $h$ while the larger trapezoid must have area $5h$. We have the equation $\frac {x}{2} = 5$. $x$ = 10, and our answer is $\boxed{\textbf{(C) }5}$

 



\textbf{Solution 4 (SIMPLE)}

By setting the area of $ABEF$ equal to $2$ times the area of $DCEF$, you get that $AB + FE = 2(DC + FE)$. Thus, $AB = 2DC + FE$. Note that $FE = DC + \frac{AB-DC}{2}$. From this, we get $AB = 2DC + DC + AB/2 - DC/2$. Simplifying, $AB/DC$ equals $\boxed{\textbf{(C) }5}$

 


	\item (\textit{AMC 10 10B 2005}) Let $x$ and $y$ be two-digit integers such that $y$ is obtained by reversing the digits
of $x$. The integers $x$ and $y$ satisfy $x^2 - y^2 = m^2$ for some positive integer $m$.
What is $x + y + m$?


 $\mathrm{(A)} 88 \qquad \mathrm{(B)} 112 \qquad \mathrm{(C)} 116 \qquad \mathrm{(D)} 144 \qquad \mathrm{(E)} 154$


 
 \textbf{Answer: }E



\textbf{Solution 1}

Let $x = 10a+b, y = 10b+a$. The given conditions imply $x>y$, which implies $a>b$, and they also imply that both $a$ and $b$ are nonzero. 

 Then, $x^2 - y^2 = (x-y)(x+y) = (9a - 9b)(11a + 11b) = 99(a-b)(a+b) = m^2$. 

 Since this must be a perfect square, all the exponents in its prime factorization must be even. $99$ factorizes into $3^2 \cdot 11$, so $11|(a-b)(a+b)$. However, the maximum value of $a-b$ is $9-1=8$, so $11|a+b$. The maximum value of $a+b$ is $9+8=17$, so $a+b=11$. 

 Then, we have $33^2(a-b) = m^2$, so $a-b$ is a perfect square, but the only perfect squares that are within our bound on $a-b$ are $1$ and $4$. We know $a+b=11$, and, for $a-b=1$, adding equations to eliminate $b$ gives us $2a=12 \Longrightarrow a=6, b=5$. Testing $a-b=4$ gives us $2a=15 \Longrightarrow a=\frac{15}{2}, b=\frac{7}{2}$, which is impossible, as $a$ and $b$ must be digits. Therefore, $(a,b) = (6,5)$, and $x+y+m=\boxed{\textbf{(E) }154}$.

 



\textbf{Solution 2}

The first steps are the same as Solution 1.  Let $x = 10a+b, y = 10b+a$, where we know that a and b are digits (whole numbers less than $10$).  

 Like Solution 1, we end up getting $(9a - 9b)(11a + 11b) = 99(a-b)(a+b) = m^2$.  This is where the solution diverges.  

 We know that the left side of the equation is a perfect square because $m$ is an integer.  If we factor $99$ into its prime factors, we get $3^2\cdot 11$.  In order to get a perfect square on the left side, $(a-b)(a+b)$ must make both prime exponents even.  Because the a and b are digits, a simple guess would be that $(a+b)$ (the bigger number) equals $11$ while $(a-b)$ is a factor of nine (1 or 9).  The correct guesses are $a = 6, b = 5$ causing $x = 65, y = 56,$ and $m = 33$.  The sum of the numbers is $\boxed{\textbf{(E) }154}$

 



\textbf{Solution 3}

Once again, the solution is quite similar as the above solutions. Since $x$ and $y$ are two digit integers, we can write $x = 10a+b, y = 10b+a$ and because $x^2 - y^2 = (x-y)(x+y)$, substituting and factoring, we get $99(a+b)(a-b) = m^2$. 

 Therefore, $(a+b)(a-b) = \frac{m^2}{99}$ and $\frac{m^2}{99}$ must be an integer. A quick strategy is to find the smallest such integer $m$ such that $\frac{m^2}{99}$ is an integer. We notice that 99 has a prime factorization of $3^2 \cdot 11.$

 Let $m^2 = n.$ Since we need a perfect square and 3 is already squared, we just need to square 11. So $3^2 \cdot 11^2$ gives us 1089 as $n$ and $m = \sqrt{1089} = 33.$ We now get the equation $(x-y)(x+y) = 33^2$, which we can also write as     $(x-y)(x+y) = 11^2 \cdot 3^2$.

 A very simple guess assumes that $x-y=3^2$ and $x+y=11^2$ since $x$ and $y$ are positive. Finally, we come to the conclusion that $x=65$ and $y=56$, so $x+y+m$ $=$ $\boxed{\textbf{(E) }154}$.

 Note that all of the solutions used $a+b$ or $a-b$ as part of their solution.

 



\textbf{Solution 4}

Continue the same as Solution $3$ until we get $33$. Knowing that $33^2 = 1089$, we have narrowed down our Pythagorean triples. We know that the $2$ other squares should be larger than $33^2$, so we can start testing. 

 If we start testing the $40$s, it is fruitless since the closest to $33^2$ would be $33 - 45 - 54$ which is not a Pythagorean triple. We can start by testing out the $50$s, and it turns our that $33 - 56 - 65$ is a Pythagorean triple. Therefore, our answer is $33+56+65$ = $\boxed{\textbf{(E) }154}$.

 



\textbf{Solution 5}

Playing around with two-digit numbers that are reverses of each other, we quickly realize they always differ by a multiple of 9 and always sum to a multiple of 11. We simplify the given expression into a difference of squares:

 
 \[(x+y)(x-y)=m^2\]

 We then take the square root of the left side of the equation

 
 \[\sqrt{(x+y)(x-y)}\]

 and notice that $(x-y)$ will always be a multiple of $9$, as we said before, which will come out of the radical right away if we assume that it is just $9$. We know that $(x+y)$ will always have a factor of $11$, so to break this part out of the radical and create an integer solution for $m$, we need another factor of 11 to make $(x+y)$ equal to 121. We simplify the equation and see that $m$ is 33. To solve for $x$ and $y$ we solve the system of $(x+y)=121$ and $(x-y)=9$ to find that $x=65$ and $y=56$. We sum the three values together and see that the answer is $\boxed{\textbf{(E) }154}$.

 


	\item (\textit{AMC 10 10B 2005}) A subset $B$ of the set of integers from $1$ to $100$, inclusive, has the property that no two elements of $B$ sum to $125$. What is the maximum possible number of elements in $B$?


 $\mathrm{(A)} 50 \qquad \mathrm{(B)} 51 \qquad \mathrm{(C)} 62 \qquad \mathrm{(D)} 65 \qquad \mathrm{(E)} 68$


 
 \textbf{Answer: }C



\textbf{Solution 1}

The question asks for the maximum possible number of elements. The integers from $1$ to $24$ can be included because you cannot make $125$ with integers from $1$ to $24$ without the other number being greater than $100$. The integers from $25$ to $100$ are left. They can be paired so the sum is $125$: $25+100$, $26+99$, $27+98$, $\ldots$, $62+63$. That is $38$ pairs, and at most one number from each pair can be included in the set. The total is $24 + 38 = \boxed{\textbf{(C)}\ 62}$.

 Also, it is possible to see that since the numbers $1$ to $24$ are in the set there are only the numbers $25$ to $100$ to consider. As $62+63$ gives $125$, the numbers $25$ to $62$ can be put in subset $B$ without having two numbers add up to $125$. In this way, subset $B$ will have the numbers $1$ to $62$, and so the answer is $\boxed{\textbf{(C)}\ 62}$.

 



\textbf{Solution 1 Alternate Solution}

Since there are $38$ numbers that sum to $125$, there are $100-38=62$ numbers not summing to $125.$~mathboy282

 



\textbf{Solution 2 (If you have no time)}

"Cut" $125$ into half. The maximum integer value in the smaller half is $62$. Thus the answer is $\boxed{\textbf{(C)}\ 62}$.

 



\textbf{Solution 3}

The maximum possible number of elements includes the smallest numbers. So, subset $B = \{1,2,3....n-1,n\}$  where n is the maximum number of elements in subset $B$. So, we have to find two consecutive numbers, $n$ and $n+1$, whose sum is $125$. Setting up our equation, we have $n+(n+1) = 2n+1 = 125$. When we solve for $n$, we get $n =\boxed{\textbf{(C)}\ 62}$.

 



\textbf{Solution 4}

We can put all odd numbers into subset B, or we can put all even numbers into subset B, so now there are 50 numbers in the set. I will use all even numbers in this solution. Now, we need to add other odd(or even!) numbers possible in this subset, which is all odd(or even) numbers that can be added so that the sum with 100(or 99) plus the biggest possible odd number(or even) to get 123. This will get us the numbers 1,3,5...21,23(or numbers 2,4,6...22,24), which gives us 12 more numbers, and adding that to the 50 original numbers, we get $B =\boxed{\textbf{(C)}\ 62}$.

 



\textbf{Solution 5 (Invariants)}

Consider $\{1,2,3,...62\}$. If we add any other integer $n$ between 1 and 100, we will have to take one out that would sum to 125 with $n$. Thus, we can change the elements, but not the number of elements (obviously, we shouldn't just take one integer out). Therefore, the maximum number of elements in this set is $|B| =\boxed{\textbf{(C)}\ 62}$.

 


\end{enumerate}